\documentclass{subfiles}
\begin{document}
    Non-singular codes are not enough. 
        More often than not, it happens that the sequence of bit in analysis may be factored
        as concatenation of distinct codewords, one again making the decoding process tedious. 

    Briefly speaking, what we are looking for are codes for which any sequence of codewords
        corresponds to at most one sequence of source symbol.
        We call such code \gls{ud} codes.
        
    Hence the questions: how do we construct UD codes?
        How do we determine if a given code is UD? 
        In the present section we answer to the latter question;
        the answer to the first question will be discussed later.

    In 1953, Sardinas and Patterson proved a result that defines an implicit algorithm 
        to determine if a code is UD. The result, stated below, is based on a simple observation:
        if we try to construct two factorizations, each word in the first factorizations
        is either a sub-word in the other factorization or has a prefix that is a suffix of 
        the words in the other factorization.

    \begin{theorem*}[Sardinas-Patterson]
        Let \(\C*\) be a code over the alphabet \(\Sigma\).
        Define the sets \(S_{0}, S_{1}, \ldots\) as follow:
        \[\begin{aligned}
            S_{0} & = \C* \\ 
            S_{i} & = \set{w \in \Sigma^{*}}[\exists \alpha \in S_{0},
                \exists \beta \in S_{i - 1} : \alpha = \beta w \lor \beta = \alpha w].
        \end{aligned}\]
        Then, necessary and sufficient condition for \(\C*\) to be UD is that,
        for all \(n > 0\), \(S_{0} \cap S_{n} = \varnothing\).
    \end{theorem*}

    Algorithmically speaking, since the number of possible sets it's finite,
        the process will eventually stop. 
        The condition in the above theorem, defines three possible halt conditions,
        namely: 
        \begin{enumerate}
            \item \(\exists i \ge 0 \text{ s.t.} S_{0} \cap S_{i} \ne \varnothing\),
                meaning the code is not UD;
            \item \(\exists i \ge 0, \text{ s.t. } S_{i} = \varnothing\),
                making the code UD; and 
            \item \(\exists i, j \ge 0 \text{ s.t. } S_{i} = S_{j}\),
                meaning the code is UD.
        \end{enumerate}

    In the example following we use colours to show the relatioship between \(\alpha, \beta\)
    and \(w\). Additionaly, we align each \(w\) to the row of the corresponfing \(\alpha\).
    \clearpage 

    \begin{example*}
        Consider the code \(\C* = \set{a, c, ad, abb, bad, deb, bbcde}\), is it UD?
            Let's apply the Sardinas-Patterson algorithm and check. 
            What we get is the following table.

            \subfile{../../TikZ/Table - Example of the Sardinas-Patterson algorithm}

            \noindent Since \(S_{0} \cap S_{5} = \set{ab} \ne \varnothing\),
            we can conclude that \(\C*\) is not UD.
    \end{example*}
\end{document}
